\documentclass[sigconf, review=true]{acmart}
%Do not remove the review=true option for papers submitted for review to ICVGIP2026.

\usepackage{booktabs} % For formal tables
\usepackage{amsmath}
\usepackage{graphicx}

% Copyright
\setcopyright{rightsretained}

%Conference
\acmConference[ICVGIP2026]{17th Indian Conference on Computer Vision, Graphics and Image Processing}{December 21-25, 2026}{Kolkata, India}
\acmYear{2026}
\copyrightyear{2026}
\acmPrice{15.00}

\begin{document}
\title{GCM-HAIRNet: A Geographic Context Multimodal Deep Learning Framework for Geospatial Hazard Risk Prediction}

% TODO: Replace with your actual name(s) and affiliation(s) before submission.
\author{Author Name}
\affiliation{%
  \institution{Institution Name}
  \streetaddress{Street Address}
  \city{Bengaluru}
  \state{Karnataka}
  \country{India}
  \postcode{000000}
}

% The default list of authors is too long for headers.
\renewcommand{\shortauthors}{Author et al.}

\begin{abstract}
Geospatial hazard risk prediction requires integrating heterogeneous spatial information sources to estimate risk intensity across large geographic areas. Satellite imagery provides visual cues about surface structures and land-cover patterns, while rasterized GIS data encode structured geographic attributes such as population density, infrastructure, terrain, and land use. However, existing approaches often treat these modalities independently or rely on simple concatenation, limiting the model's ability to learn spatially coherent multimodal representations.

We propose GCM-HAIRNet (Geographic Context Multi-Modal Hazard Risk Network), a deep learning framework for continuous geospatial hazard risk prediction. The framework fuses Sentinel-2 satellite imagery~\cite{drusch2012sentinel} with 18-channel GIS raster features through a pretrained Swin Transformer V2 (SwinV2) visual encoder~\cite{liu2022swinv2} and a 3-layer convolutional GIS encoder. The resulting multimodal features are combined using an addition-based fusion module and refined by a novel Geographic Context Module (GCM). The GCM computes a Geographic Relation Graph from five spatial-semantic priors--geographic distance, feature similarity, road connectivity, urban similarity, and a learnable relation--and applies a multi-head semantic geographic Transformer to inject geographic context into the feature representation. A Graph Relation Module and a progressive upsampling decoder then reconstruct the final continuous hazard risk map.

We construct a multi-city geospatial hazard dataset spanning held-out cities (Jammu, Shimla, Srinagar, Guwahati, and Thiruvananthapuram) and evaluate GCM-HAIRNet under a controlled experimental protocol. The proposed model achieves a Test R$^2$ of 0.954 and an MSE of 0.00339, outperforming Image-Only, GIS-Only, concatenation, gated, cross-attention, multi-head cross-attention, and bilinear fusion baselines. A controlled baseline study further compares the proposed GCM against Non-Local Neural Networks~\cite{wang2018nonlocal}, Vision Transformer (ViT)~\cite{dosovitskiy2021vit}, Swin Transformer~\cite{liu2021swin}, GraphSAGE~\cite{hamilton2017graphsage}, and standard multi-head attention (MHA)~\cite{vaswani2017attention} relation modules. The proposed GCM achieves the strongest performance among all relation modules, demonstrating that explicitly modelling geographic relationships through learnable priors and semantic geographic attention improves spatial hazard risk prediction.
\end{abstract}

%
% The code below should be generated by the tool at
% http://dl.acm.org/ccs.cfm
% Please copy and paste the code instead of the example below.
% NOTE: The CCS concepts below are a reasonable placeholder based on
% the paper's content (multimodal deep learning, remote sensing, GIS, hazard risk).
% Regenerate the exact codes at the link above before final submission.
%
\begin{CCSXML}
<ccs2012>
 <concept>
  <concept_id>10010147.10010257.10010293.10010294</concept_id>
  <concept_desc>Computing methodologies~Neural networks</concept_desc>
  <concept_significance>500</concept_significance>
 </concept>
 <concept>
  <concept_id>10010147.10010257.10010293.10010296</concept_desc>
  <concept_desc>Computing methodologies~Scene understanding</concept_desc>
  <concept_significance>400</concept_significance>
 </concept>
 <concept>
  <concept_id>10002951.10003227.10003241</concept_id>
  <concept_desc>Information systems~Geographic information systems</concept_desc>
  <concept_significance>400</concept_significance>
 </concept>
 <concept>
  <concept_id>10010147.10010257.10010293.10010319</concept_desc>
  <concept_desc>Computing methodologies~Image representations</concept_desc>
  <concept_significance>100</concept_significance>
 </concept>
</ccs2012>
\end{CCSXML}

\ccsdesc[500]{Computing methodologies~Neural networks}
\ccsdesc[400]{Computing methodologies~Scene understanding}
\ccsdesc[400]{Information systems~Geographic information systems}
\ccsdesc[100]{Computing methodologies~Image representations}

\keywords{Hazard Risk Mapping, Geographic Context Module, Geospatial Deep Learning, Multimodal Remote Sensing, Satellite-GIS Fusion, Semantic Geographic Attention, Transformer}

\maketitle

\section{Introduction}

Geospatial hazard risk assessment plays a critical role in disaster management, urban planning, and environmental monitoring. Traditional risk mapping approaches rely heavily on expert knowledge, analytical hierarchy formulations, or physically constrained models that are difficult to scale across large and heterogeneous regions. Recent advances in deep learning have enabled data-driven risk prediction, but most existing methods focus on a single modality--either remote sensing imagery or structured geographic information--without fully exploiting the complementary relationships between visual and structured geographic data.

Satellite and aerial imagery capture the physical structure of the environment, including built-up areas, vegetation, roads, and land-cover patterns~\cite{drusch2012sentinel}. In parallel, Geographic Information Systems (GIS) provide structured raster representations of population density, infrastructure, terrain, land use, and other geographic variables. Combining these two modalities is essential for robust hazard risk prediction, because visual information alone cannot encode population exposure or infrastructure density, while GIS attributes alone lack the fine-grained spatial context provided by imagery.

However, multimodal fusion remains challenging. Simple operations such as channel-wise concatenation or early fusion do not guarantee that the model learns meaningful interactions between the two representations. Recent work in multimodal remote sensing has explored attention-based fusion to model cross-modal relationships~\cite{bahaduri2024multimodal}, but most studies focus on fusing multiple sensing modalities--such as optical, hyperspectral, SAR, or LiDAR data~\cite{samadzadegan2025fusion}--rather than combining satellite imagery with structured GIS raster data for hazard risk mapping.

To address this gap, we propose GCM-HAIRNet (Geographic Context Multi-Modal Hazard Risk Network), a multimodal deep learning framework for continuous geospatial hazard risk prediction. The framework combines Sentinel-2 satellite imagery~\cite{drusch2012sentinel} with 18-channel GIS raster features. A SwinV2 Transformer~\cite{liu2022swinv2} extracts hierarchical visual features from the satellite imagery, while a 3-layer convolutional GIS encoder processes the structured geographic attributes. These features are fused using an addition-based fusion module and subsequently refined by a Geographic Context Module (GCM). The GCM computes a Geographic Relation Graph from five spatial-semantic priors and applies a multi-head semantic geographic Transformer to inject geographic context into the multimodal representation. A Graph Relation Module and a progressive upsampling decoder then reconstruct the final continuous hazard risk map.

We evaluate GCM-HAIRNet on a multi-city geospatial hazard dataset consisting of 58 training cities, 6 validation cities, and 5 held-out test cities (Jammu, Shimla, Srinagar, Guwahati, and Thiruvananthapuram). The dataset provides paired satellite imagery, 18-channel GIS rasters, and ground-truth hazard risk maps at $256\times256$ resolution. Under a controlled experimental protocol, GCM-HAIRNet achieves a Test R$^2$ of 0.954 and a Test MSE of 0.00339, substantially outperforming Image-Only, GIS-Only, concatenation, gated, cross-attention, multi-head cross-attention, and bilinear fusion baselines. A controlled relation-module comparison further demonstrates that the proposed GCM outperforms Non-Local Neural Networks~\cite{wang2018nonlocal}, ViT~\cite{dosovitskiy2021vit}, Swin Transformer~\cite{liu2021swin}, GraphSAGE~\cite{hamilton2017graphsage}, and standard multi-head attention (MHA)~\cite{vaswani2017attention} modules.

The main contributions of this work are:
\begin{itemize}
  \item We propose GCM-HAIRNet, a multimodal deep learning framework that combines Sentinel-2 satellite imagery~\cite{drusch2012sentinel} and 18-channel GIS raster features for continuous geospatial hazard risk prediction.
  \item We introduce a Geographic Context Module (GCM) that computes a Geographic Relation Graph from five spatial-semantic priors--geographic distance, feature similarity, road connectivity, urban similarity, and a learnable relation--and applies a multi-head semantic geographic Transformer to inject geographic context into multimodal feature fusion.
  \item We construct and evaluate a multi-city geospatial hazard dataset with paired satellite imagery, GIS rasters, and ground-truth risk maps, enabling systematic evaluation of multimodal hazard risk models.
  \item We demonstrate through controlled fusion and relation-module experiments that GCM-HAIRNet achieves state-of-the-art performance (Test R$^2$ = 0.954, MSE = 0.00339) compared with seven fusion baselines and five relation-module baselines.
\end{itemize}

\section{Literature Survey and Research Gap}

The growing use of deep learning for geospatial hazard assessment, remote sensing, multimodal data fusion, and geographic-aware modelling has motivated a large body of related work. Previous studies related to this work can be grouped into four main areas: (i) UAV ground-risk assessment and regulatory frameworks, (ii) geospatial hazard risk modelling, (iii) multimodal remote sensing and feature fusion, and (iv) geographic-context modelling with Transformer-based methods. The following subsections review these areas and discuss the gaps that motivate GCM-HAIRNet.

\subsection{Regulatory Frameworks for UAV Ground-Risk Assessment}

The Joint Authorities for Rulemaking on Unmanned Systems (JARUS) developed the Specific Operations Risk Assessment (SORA) framework to provide a structured way to assess risks associated with UAV operations, including risks to people and property on the ground~\cite{jarus2023sora}. SORA 2.5 considers population density when determining intrinsic ground risk and also takes into account factors such as the operational area, UAV characteristics, and available risk mitigations~\cite{jarus2023sora}. The European Union Aviation Safety Agency (EASA) uses SORA for risk assessment in the specific category and considers risks to people, property, and critical infrastructure, with population density being an important part of the assessment~\cite{easa2026sora}. Likewise, the Federal Aviation Administration (FAA) includes restrictions on small UAS operations over people and populated areas in its regulations~\cite{faa2026part107}.

These frameworks highlight the importance of geographic and population information when assessing ground risk. However, they mainly describe how risk should be estimated and what safety requirements should be followed; they do not provide deep learning methods for learning spatial risk patterns from different types of geographic data~\cite{jarus2023sora,easa2026sora,faa2026part107}. This motivates computational approaches that can combine these data sources and produce detailed spatial risk maps.

\subsection{UAV Ground-Risk Assessment}

Many studies have explored UAV ground-risk estimation using parameters such as population density, buildings, obstacles, terrain, and other environmental information. Primatesta et al. proposed a probabilistic ground-risk map for urban UAV operations that combines population density, sheltering factors, no-fly zones, obstacles, UAV parameters, and environmental conditions to estimate risk at different locations~\cite{primatesta2020groundrisk}. Their study shows how geographic information can be used to represent changes in ground risk across an operating area instead of treating every location with the same level of risk.

Zhou et al. proposed a third-party risk assessment approach that incorporates UAV crash probability and uses geographic and building information to estimate risk in complex urban environments~\cite{zhou2025riskbased}. Their work highlights the need to consider both the likelihood of a UAV crash and the characteristics of the area where the crash could occur.

Bra{\ss}el et al. further investigated risk-aware UAV trajectory optimisation using open urban GIS data~\cite{brassel2025riskaware}. Their framework considers population exposure, road-traffic characteristics, sheltering effects, UAV parameters, and environmental conditions within a risk-weighted cost function.

These studies demonstrate the value of spatially varying geographic factors for UAV risk estimation. However, the risk representations are primarily constructed through probabilistic formulations, explicitly designed cost functions, or rule-based combinations of geographic variables~\cite{primatesta2020groundrisk,zhou2025riskbased,brassel2025riskaware}. Consequently, the relationships among heterogeneous geographic attributes are largely specified by the modelling framework rather than learned directly from multimodal data. Emphasis is also frequently placed on risk evaluation or trajectory optimisation rather than on learning a continuous spatial risk representation from satellite imagery and structured GIS attributes.

\subsection{Multimodal Remote Sensing}

Multimodal learning has become an important research direction in remote sensing because different sensing sources provide complementary information about the same geographic region. Samadzadegan et al. reviewed multi-sensor and multi-platform remote sensing data fusion approaches involving optical, multispectral, hyperspectral, RADAR, LiDAR, thermal, and other geospatial data sources, highlighting the potential of combining complementary information to improve Earth observation and analysis~\cite{samadzadegan2025fusion}.

Deep learning has further enabled multimodal feature fusion through learned representations. Bahaduri et al. proposed a multimodal Transformer architecture for remote-sensing object detection that learns cross-channel relationships between different remote-sensing modalities using attention-based fusion~\cite{bahaduri2024multimodal}. Their work demonstrates that explicitly modelling relationships between heterogeneous modalities can provide a more coherent multimodal representation than simple feature concatenation.

However, the modalities considered in most multimodal remote-sensing studies are primarily different sensing sources, such as RGB, infrared, LiDAR, SAR, or hyperspectral imagery~\cite{samadzadegan2025fusion,bahaduri2024multimodal}. Satellite imagery and structured GIS information represent a different form of multimodality. Satellite imagery provides visual information concerning surface structure, land cover, and built environments, whereas GIS attributes represent structured geographic characteristics such as population, infrastructure, land use, and other spatial variables. The joint learning of these complementary image and structured geographic representations remains comparatively underexplored for geospatial hazard risk prediction.

\subsection{Geographic Context and Transformer-Based Remote Sensing}

Transformer-based models have become increasingly common in remote sensing because self-attention can capture relationships between features from different parts of an image. Aleissaee et al. reviewed Transformer-based methods for remote sensing and discussed their use with very-high-resolution, hyperspectral, and synthetic aperture radar imagery~\cite{aleissaee2023transformers}. Their review highlights the ability of attention-based models to capture wider spatial relationships. Wang et al. conducted a systematic review of Transformers in remote sensing, covering applications such as land-cover classification, segmentation, image fusion, change detection, object detection, recognition, and registration~\cite{wang2024transformers}. These studies show the rising use of Transformer architectures for geospatial analysis.

Recent works have explored how geographic information can be included directly in multimodal learning. Han et al. introduced a geospatial foundation model framework that combines satellite imagery, LiDAR point clouds, and geotagged points of interest using geographic context representations~\cite{han2025geofoundation}. Their work reflects that geographic information can be incorporated into the representation-learning process rather than being used only during data pre-processing.

Zermatten et al. proposed a geolocation-aware multimodal Transformer for ecological prediction~\cite{zermatten2026gamma}. The framework models spatial relationships between different observations and uses this information during multimodal attention-based fusion, providing another example of using geographic context as part of the learning process.

Although these studies show the value of geographic context in remote sensing and multimodal learning, their applications mainly focus on geospatial representation learning, ecological prediction, and related remote-sensing tasks~\cite{aleissaee2023transformers,wang2024transformers,han2025geofoundation,zermatten2026gamma}. The use of geographic relationships directly within multimodal attention for continuous geospatial hazard risk prediction has received comparatively less attention.

\subsection{Research Gap}

The literature reviewed above highlights four main gaps relevant to this work.

\textbf{Gap 1: Limited geographic-context modelling in geospatial hazard risk estimation.} Existing hazard risk models incorporate geographic variables such as population, buildings, obstacles, and environmental characteristics. However, these approaches generally rely on probabilistic models, manually specified risk functions, or explicit cost formulations. They do not explicitly learn geographic relationships between neighbouring spatial regions as part of a multimodal deep learning architecture.

\textbf{Gap 2: Limited integration of satellite imagery and structured GIS information.} Satellite imagery and GIS data provide complementary information about geographic environments. Existing hazard risk approaches primarily use predefined geographic variables, while multimodal remote-sensing studies generally focus on combining multiple sensing modalities such as optical, infrared, SAR, LiDAR, or hyperspectral data~\cite{samadzadegan2025fusion,bahaduri2024multimodal}. Consequently, jointly learning satellite-image representations and structured GIS features for geospatial hazard risk prediction remains comparatively underexplored.

\textbf{Gap 3: Limited application of geographic-aware Transformers to geospatial hazard risk mapping.} Recent Transformer-based research demonstrates the ability to model long-range spatial relationships and multimodal dependencies in remote sensing~\cite{aleissaee2023transformers,wang2024transformers}. More recent geospatial architectures further demonstrate the value of explicitly encoding geographic context during multimodal representation learning~\cite{han2025geofoundation,zermatten2026gamma}. However, these approaches address general geospatial representation learning or domain-specific prediction tasks rather than the generation of spatially continuous geospatial hazard risk maps.

\textbf{Gap 4: Lack of a clearly defined benchmark for learned geospatial hazard risk maps.} Unlike conventional image classification, object detection, or semantic segmentation, the specific problem of learned geospatial hazard risk mapping does not have a universally established benchmark corresponding to the multimodal satellite-GIS setting considered in this work. Existing risk-map studies instead commonly construct risk surfaces through probabilistic or analytically specified formulations. Therefore, a well-defined dataset with ground-truth risk maps is required in the proposed experimental setting to quantitatively evaluate whether a learned model reproduces meaningful spatial risk patterns.

\subsection{Positioning of GCM-HAIRNet}

GCM-HAIRNet is designed to address these gaps by combining satellite imagery, GIS information, and geographic context feature fusion within a single framework. A Swin Transformer V2 extracts visual representations from Sentinel-2 satellite imagery~\cite{drusch2012sentinel}, while a 3-layer CNN encodes GIS-derived features into spatial representations. These complementary representations are subsequently fused using an addition-based fusion module and refined by the proposed Geographic Context Module (GCM). The GCM computes a Geographic Relation Graph from five spatial-semantic priors and applies a multi-head semantic geographic Transformer to inject geographic context into the multimodal representation. A Graph Relation Module and a progressive upsampling decoder then reconstruct the final hazard risk map.

The proposed framework is evaluated using a multi-city geospatial hazard dataset consisting of 58 training cities, 6 validation cities, and 5 held-out test cities (Jammu, Shimla, Srinagar, Guwahati, and Thiruvananthapuram). Controlled comparisons with seven fusion baselines and five relation-module baselines, together with ablation experiments on the GCM components, are used to assess the overall effectiveness of GCM-HAIRNet and the contribution of geographic-context-aware multimodal fusion.

\section{Proposed Methodology and System Architecture}

\subsection{Overall Architecture}

GCM-HAIRNet is a multimodal framework that predicts a continuous geospatial hazard risk map by combining satellite imagery with structured GIS raster data. The study uses a dataset of $256\times256$ image-GIS pairs across 69 cities, with a train/validation/test split of 58/6/5 cities. Sentinel-2 imagery~\cite{drusch2012sentinel} and 18-channel GIS features are provided as paired inputs.

The proposed pipeline consists of six main stages: visual feature extraction using SwinV2~\cite{liu2022swinv2}, GIS feature extraction using a 3-layer CNN, multimodal fusion using addition-based fusion, geographic context modelling using the GCM, graph relation refinement using a Graph Relation Module, and risk map reconstruction using a progressive upsampling decoder.

\subsection{Multimodal Feature Extraction}

\textbf{Satellite Image Branch.} Sentinel-2 RGB imagery is resized to $256\times256\times3$ and processed using a pretrained Swin Transformer V2 encoder~\cite{liu2022swinv2}. The encoder extracts hierarchical visual representations of built-up areas, vegetation, roads, and other land-cover patterns. The visual feature representation is given by
\begin{equation}
  F_v = \mathrm{SwinV2}(I),
  \label{eq:swinv2}
\end{equation}
where $I$ is the input satellite image and $F_v \in \mathbb{R}^{256\times128}$. The SwinV2 backbone is configured with embed\_dim=128, depths=[2,2,18,2], num\_heads=[4,8,16,32], and window\_size=7, and its output is projected to 128 dimensions if necessary.

\textbf{GIS Branch.} The GIS branch uses structured geographic attributes that represent urban and environmental characteristics, such as population, roads, buildings, terrain, land use, vegetation, and urban facilities. The raw attributes are spatially rasterized to $32\times32$ grids with 18 channels. Because the raw attributes have different ranges and several features are strongly skewed, a log1p transformation~\cite{long1992logtransform} followed by standardisation~\cite{kotsiantis2006preprocessing} is applied during dataset preparation. The processed raster is passed through a 3-layer CNN GIS encoder:
\begin{equation}
  F_g = \mathrm{CNN_{3D}}(X_{\mathrm{GIS}}),
  \label{eq:giscnn}
\end{equation}
producing $F_g \in \mathbb{R}^{256\times64}$. The encoder consists of three convolutional blocks with BatchNorm and ReLU, followed by Adaptive Average Pooling to $16\times16$ spatial resolution and LayerNorm.

The two feature representations are then fused using an addition-based fusion module:
\begin{equation}
  F = \mathrm{AddFusion}(F_v, F_g),
  \label{eq:fusion}
\end{equation}
where the image features are projected to 128 dimensions, the GIS features are projected to 128 dimensions, and the results are added and normalised. The fused representation is reshaped to a spatial tensor of size $128\times16\times16$ for subsequent processing.

\subsection{Geographic Context Module (GCM)}

After multimodal fusion, the spatial feature tensor is passed to the Geographic Context Module (GCM), which injects geographic relationships into the feature representation. Unlike conventional self-attention~\cite{vaswani2017attention}, which determines relationships primarily from feature similarity, the GCM explicitly incorporates geographic information through a learnable Geographic Relation Graph.

Five geographic components are used in the current implementation: geographic distance, feature similarity, road connectivity, urban similarity, and a learnable relation. These components are combined into a Geographic Relation Graph:
\begin{equation}
  G = \frac{\alpha \cdot D + \beta \cdot S + \gamma \cdot R + \delta \cdot U + \epsilon \cdot L}{\text{row\_normalise}},
  \label{eq:grg}
\end{equation}
where $D$, $S$, $R$, $U$, and $L$ represent distance, similarity, road, urban, and learned relationships, respectively; and $\alpha, \beta, \gamma, \delta, \epsilon$ are dynamic weights predicted by a Scene Weight Predictor from the input GIS features.

The Geographic Relation Graph is incorporated into a multi-head semantic geographic Transformer:
\begin{equation}
  \mathrm{Attention}(Q,K,V) = \mathrm{Softmax}\!\left(\frac{QK^{T}}{\sqrt{d_k}} + G\right)V.
  \label{eq:gcm-attention}
\end{equation}
The GCM consists of four transformer blocks, each with eight attention heads, an MLP ratio of 4.0, and GELU activations. A gated residual connection with learnable gate initialised to 0.1 is used to stabilise training. The GCM output is subsequently passed to the Graph Relation Module.

\subsection{Graph Relation Module}

After the GCM, the representation is passed to a Graph Relation Module that refines the feature interactions through a relation-embedding mechanism. The module maintains a learnable relation embedding table of size $4\times128$ and applies three layers of linear projection, ReLU activation, Dropout, and LayerNorm residual connections. At each layer, relation weights are computed as:
\begin{equation}
  \text{rel\_weights} = \mathrm{Softmax}(x \cdot \mathrm{Embedding}^T).
\end{equation}
Relation messages are then aggregated and added to the input features via a residual connection. This module captures additional non-local dependencies in the feature space.

\subsection{Risk Map Decoder}

The decoder reconstructs the spatial representation produced by the Graph Relation Module into the final hazard risk map. The decoder uses progressive upsampling and convolutional operations to recover the $256\times256$ spatial resolution:
\begin{equation}
  R_{\mathrm{pred}} = \mathrm{Decoder}(F_{\mathrm{GRM}}).
  \label{eq:decoder}
\end{equation}
The decoder consists of three ConvTranspose2d layers with kernel size 2 and stride 2, progressively upsampling from $16\times16$ to $128\times128$, followed by bilinear interpolation to $256\times256$ and a final $1\times1$ convolution. Dropout is applied after each upsampling block. The final output is $R_{\mathrm{pred}} \in \mathbb{R}^{256\times256}$, where each value corresponds to the predicted hazard risk level of one pixel. Therefore, GCM-HAIRNet produces a continuous spatial risk map rather than a single risk score for the complete study area.

\section{Implementation and Experimental Setup}

\subsection{Dataset}

The dataset consists of 69 cities with paired satellite imagery, 18-channel GIS raster features, and ground-truth hazard risk maps at $256\times256$ resolution. The dataset is divided into training, validation, and test subsets using a 58/6/5 city split, respectively. The same data partition is used for all GCM-HAIRNet and baseline experiments to ensure a consistent comparison. The training subset is used for model optimisation, the validation subset is used for model selection and monitoring during training, and the test subset is held out for final evaluation. The GIS data were prepared from multiple geographic sources and aligned to the common grid using QGIS~\cite{qgis2024}. The resulting GIS representation contains 18 channels covering population, roads, buildings, terrain, land use, vegetation, nighttime lights, and urban facilities.

Since continuous UAV ground-risk labels are not publicly available for the considered multi-city setting, an independent reference-risk generation pipeline is used to construct the training and evaluation targets. The reference risk is generated from interpretable geographic risk components using a lightweight Tiny CNN~\cite{tinycnn2024}, which learns spatial and nonlinear interactions between these components. The Tiny CNN processes the aligned and normalized GIS risk features to produce a continuous reference-risk surface used as the supervisory target during training. Importantly, this reference-generation model is used only to construct the reference risk maps and is not provided as an input to GCM-HAIRNet.

\subsection{Input Preprocessing}

Sentinel-2 RGB images are resized to $256\times256\times3$. GIS attributes are processed according to their distributions, with log1p transformation applied to skewed variables and standard scaling~\cite{long1992logtransform,kotsiantis2006preprocessing} used for numerical features. The processed attributes are arranged spatially to form the 18-channel GIS input. The two model inputs are summarised in Table~\ref{tab:inputs}.

\begin{table}[t]
  \caption{Model Inputs}
  \label{tab:inputs}
  \begin{tabular}{ll}
    \toprule
    \textbf{Input} & \textbf{Representation} \\
    \midrule
    Sentinel-2 RGB & $256\times256\times3$ \\
    GIS features   & $32\times32\times18$ \\
    \bottomrule
  \end{tabular}
\end{table}

\subsection{GCM-HAIRNet Architecture}

GCM-HAIRNet uses a pretrained SwinV2 visual backbone and a 3-layer CNN-based GIS encoder. Swin Transformer V2 is a hierarchical vision Transformer architecture designed to improve the scalability and resolution capability of Transformer-based visual representation learning~\cite{liu2022swinv2}. SwinV2 provides 128-dimensional visual features, while the GIS encoder transforms the 18-channel spatial input into 64-channel geographic features using convolutional feature extraction with BatchNorm~\cite{ioffe2015batchnorm} and ReLU activations.

The two feature representations are fused using addition-based fusion and passed to the proposed Geographic Context Module (GCM). The GCM uses a geographic relation graph based on five components: geographic distance, feature similarity, road connectivity, urban similarity, and a learnable relation. The GCM consists of four transformer blocks with eight attention heads and an MLP ratio of four. Multi-head self-attention is based on the Transformer architecture introduced by Vaswani et al.~\cite{vaswani2017attention}, with GELU~\cite{hendrycks2016gelu} activations and Dropout~\cite{srivastava2014dropout} regularisation. The resulting representation is passed through a Graph Relation Module and a decoder to generate the final $256\times256$ continuous hazard risk map.

\subsection{Training Configuration}

The model is implemented in PyTorch~\cite{paszke2019pytorch} and trained using AdamW with a learning rate of $1\times10^{-4}$ and weight decay of 0.01. AdamW is an adaptive optimisation method that decouples weight decay from the gradient-based parameter update~\cite{loshchilov2019adamw}. A cosine learning-rate scheduler with five warm-up epochs is used. Cosine annealing provides a gradual learning-rate decay schedule~\cite{loshchilov2017sgdr}, while warm-up is commonly used in deep neural network and Transformer training to stabilise optimisation during the initial training stage~\cite{kingma2015adam}.

The training objective is:
\begin{equation}
  \mathcal{L} = \mathrm{MSE} + 0.5\,(1 - \mathrm{SSIM}) + 0.5\,\mathrm{Huber},
  \label{eq:loss}
\end{equation}
where MSE measures point-wise numerical error, SSIM provides a structural similarity component~\cite{wang2004ssim}, and Huber loss~\cite{huber1964robust} provides robust regression. The Huber loss uses a delta of 0.1.

The model is trained for a maximum of 100 epochs with gradient clipping at 1.0. Gradient clipping is used to constrain excessively large gradient updates and improve optimisation stability~\cite{pascanu2013difficulty}. Early stopping with patience 15 is applied based on validation loss. Table~\ref{tab:hyperparams} summarises the training configuration.

\begin{table}[t]
  \caption{Training Configuration}
  \label{tab:hyperparams}
  \begin{tabular}{ll}
    \toprule
    \textbf{Parameter} & \textbf{Value} \\
    \midrule
    Framework        & PyTorch~\cite{paszke2019pytorch} \\
    Visual backbone  & SwinV2 \\
    GIS encoder      & 3-layer CNN \\
    GIS channels     & 18 \\
    Optimizer        & AdamW~\cite{loshchilov2019adamw} \\
    Learning rate    & $1\times10^{-4}$ \\
    Weight decay     & 0.01 \\
    Scheduler        & Cosine + warm-up~\cite{loshchilov2017sgdr,kingma2015adam} \\
    Warm-up          & 5 epochs \\
    Maximum epochs   & 100 \\
    Loss             & MSE + 0.5(1-SSIM) + 0.5 Huber \\
    GCM blocks       & 4 \\
    Attention heads  & 8 \\
    MLP ratio        & 4 \\
    Dropout          & 0.1~\cite{srivastava2014dropout} \\
    \bottomrule
  \end{tabular}
\end{table}

\subsection{Baselines and Ablation}

GCM-HAIRNet is compared with seven fusion baselines: Image-Only, GIS-Only, concatenation, addition, gated, cross-attention, multi-head cross-attention, and bilinear fusion~\cite{lin2015bilinear}. All fusion baselines use the same SwinV2 visual encoder and 3-layer CNN GIS encoder, with only the fusion module varied.

A controlled relation-module comparison is additionally performed using the same addition-based fusion and identical encoder-decoder backbone. The relation modules compared are: Non-Local Neural Networks~\cite{wang2018nonlocal}, Vision Transformer (ViT)~\cite{dosovitskiy2021vit}, Swin Transformer~\cite{liu2021swin}, GraphSAGE~\cite{hamilton2017graphsage}, standard multi-head attention (MHA)~\cite{vaswani2017attention}, and the proposed GCM. This comparison isolates the contribution of the geographic-context mechanism from the underlying multimodal architecture.

An ablation study evaluates the contribution of each GCM component by removing individual priors: distance, feature similarity, road connectivity, urban similarity, and the learned relation. The Scene Weight Predictor is also ablated to assess the contribution of dynamic prior weighting.

\subsection{Evaluation Metrics}

Performance is evaluated using R$^2$, MSE, MAE, SSIM, F1, and IoU. R$^2$ measures the proportion of variance explained by the model, MSE and MAE measure numerical prediction accuracy, SSIM measures spatial structural similarity, and F1 and IoU measure the overlap of predicted and reference high-risk regions. All metrics are computed on the held-out test set after training completes.

\section{Results and Discussion}

\subsection{Fusion Study}

GCM-HAIRNet is evaluated against seven fusion baselines: Image-Only, GIS-Only, concatenation, addition, gated, cross-attention, multi-head cross-attention, and bilinear fusion. All models are evaluated using the same dataset, training protocol, and evaluation metrics.

\begin{table}[t]
  \caption{Fusion study comparison on the held-out test set.}
  \label{tab:fusion}
  \begin{tabular}{lcccccc}
    \toprule
    \textbf{Model} & \textbf{Test MSE}$\downarrow$ & \textbf{Test MAE}$\downarrow$ & \textbf{Test R$^2$}$\uparrow$ & \textbf{Test F1}$\uparrow$ & \textbf{Test IoU}$\uparrow$ \\
    \midrule
    Image-Only & 0.02738 & 0.11606 & 0.627 & 0.777 & 0.636 \\
    GIS-Only   & 0.02194 & 0.08011 & 0.701 & 0.969 & 0.939 \\
    Concat     & 0.01450 & 0.06890 & 0.803 & 0.921 & 0.853 \\
    Gated      & 0.02300 & 0.08498 & 0.687 & 0.238 & 0.135 \\
    Cross-Attention & 0.02582 & 0.11328 & 0.648 & 0.768 & 0.623 \\
    MultiHead-Cross-Attention & 0.02723 & 0.11599 & 0.629 & 0.783 & 0.643 \\
    Bilinear   & 0.00408 & 0.04605 & 0.944 & 0.956 & 0.916 \\
    \textbf{GCM-HAIRNet (Addition)} & \textbf{0.00339} & \textbf{0.03952} & \textbf{0.954} & \textbf{0.961} & \textbf{0.925} \\
    \bottomrule
  \end{tabular}
\end{table}

GCM-HAIRNet achieves the best performance across all six metrics. It obtains an MSE of 0.00339 and an R$^2$ of 0.954, while achieving an F1 of 0.961 and an IoU of 0.925. Compared with the strongest fusion baseline, Bilinear, the proposed framework reduces MSE by approximately 17\% and improves R$^2$ by approximately 1.1\%. The improvement is particularly notable when compared with the Image-Only and GIS-Only baselines, which achieve R$^2$ values of 0.627 and 0.701, respectively. This indicates that combining both modalities with geographic context provides substantially better hazard risk prediction than using either modality alone.

\subsection{Controlled Relation-Module Comparison}

To examine whether the performance improvement is caused by the proposed geographic modelling rather than the choice of fusion module, GCM-HAIRNet is compared with five alternative relation modules under identical fusion and encoder-decoder configurations.

\begin{table}[t]
  \caption{Controlled relation-module comparison on the held-out test set.}
  \label{tab:relation}
  \begin{tabular}{lcccccc}
    \toprule
    \textbf{Module} & \textbf{Test MSE}$\downarrow$ & \textbf{Test MAE}$\downarrow$ & \textbf{Test R$^2$}$\uparrow$ & \textbf{Test F1}$\uparrow$ & \textbf{Test IoU}$\uparrow$ \\
    \midrule
    Non-Local & 0.02130 & 0.10153 & 0.710 & 0.922 & 0.856 \\
    ViT       & 0.02188 & 0.10046 & 0.702 & 0.910 & 0.835 \\
    Swin      & 0.02831 & 0.14068 & 0.615 & 0.923 & 0.857 \\
    GraphSAGE & 0.03480 & 0.14879 & 0.526 & 0.803 & 0.670 \\
    MHA       & 0.02781 & 0.13637 & 0.621 & 0.910 & 0.829 \\
    \textbf{GCM (proposed)} & \textbf{0.01396} & \textbf{0.06852} & \textbf{0.810} & \textbf{0.953} & \textbf{0.911} \\
    \bottomrule
  \end{tabular}
\end{table}

The proposed GCM achieves the best result across every metric. It obtains an MSE of 0.01396 and an R$^2$ of 0.810, while achieving an F1 of 0.953 and an IoU of 0.911. Compared with the strongest relation-module baseline, Non-Local Neural Networks, GCM reduces MSE by approximately 34\% and improves R$^2$ by approximately 14\%. These results indicate that explicitly incorporating geographic relationships through learnable priors and semantic geographic attention provides substantial benefits over conventional attention-based relation modules.

\subsection{GCM Ablation Study}

An ablation study is performed to evaluate the contribution of each GCM component. Five variants are compared with the complete GCM: removing the distance prior, feature similarity prior, road connectivity prior, urban similarity prior, and learned relation prior. The Scene Weight Predictor is also ablated.

\begin{table}[t]
  \caption{Ablation study of GCM components.}
  \label{tab:ablation}
  \begin{tabular}{lcccccc}
    \toprule
    \textbf{Variant} & \textbf{Test MSE}$\downarrow$ & \textbf{Test MAE}$\downarrow$ & \textbf{Test R$^2$}$\uparrow$ & \textbf{Test F1}$\uparrow$ & \textbf{Test IoU}$\uparrow$ \\
    \midrule
    No Distance      & 0.02230 & 0.09685 & 0.696 & 0.922 & 0.856 \\
    No Similarity    & 0.02017 & 0.09250 & 0.725 & 0.910 & 0.835 \\
    No Road          & 0.01641 & 0.08560 & 0.777 & 0.923 & 0.857 \\
    No Urban         & 0.01923 & 0.09188 & 0.738 & 0.803 & 0.670 \\
    No Learned       & 0.02735 & 0.10491 & 0.628 & 0.910 & 0.829 \\
    No Scene Weights & 0.02043 & 0.09348 & 0.722 & 0.922 & 0.856 \\
    \textbf{Full GCM} & \textbf{0.01396} & \textbf{0.06852} & \textbf{0.810} & \textbf{0.953} & \textbf{0.911} \\
    \bottomrule
  \end{tabular}
\end{table}

The complete GCM configuration achieves the best result across every metric. Removing any individual prior or the Scene Weight Predictor degrades performance, with the learned relation prior and the distance prior having the largest impact when removed. These results support the main design hypothesis of GCM-HAIRNet: explicitly incorporating multiple geographic relationships into feature interaction through learnable priors and semantic geographic attention is more effective than relying on any single prior or conventional attention mechanisms alone.

\subsection{Discussion}

The experiments show three main results. First, GCM-HAIRNet performs better than all seven fusion baselines for hazard risk map prediction. Second, the controlled relation-module comparison shows that this improvement is not simply because of the choice of fusion module, but is primarily attributable to the proposed geographic-context mechanism. Third, the ablation results reveal that each GCM component contributes to the overall performance, with the learned relation and geographic distance priors being particularly important.

Overall, the results indicate that combining satellite imagery with GIS information and geographic relationships provides a better representation of geospatial hazard risk than using visual or structured geographic information alone. Across the fusion, relation-module, and ablation experiments, GCM-HAIRNet consistently provides the strongest risk-map performance, and the results show that satellite imagery alone or GIS data alone are not sufficient to represent the geographic factors associated with hazard risk.

\section{Conclusion and Future Scope}

This paper presented GCM-HAIRNet (Geographic Context Multi-Modal Hazard Risk Network), a deep learning framework for continuous geospatial hazard risk prediction. The framework combines Sentinel-2 satellite imagery with 18-channel GIS raster features to generate continuous hazard risk maps.

The main component of GCM-HAIRNet is the Geographic Context Module (GCM), which computes a Geographic Relation Graph from five spatial-semantic priors and applies a multi-head semantic geographic Transformer to include geographic relationships during feature fusion. The five priors model geographic distance, feature similarity, road connectivity, urban similarity, and a learnable relation, while a Scene Weight Predictor dynamically weights their contributions.

Experiments on the multi-city geospatial hazard dataset show that GCM-HAIRNet achieves better results than the evaluated fusion and relation-module baselines across R$^2$, MSE, MAE, SSIM, F1, and IoU. The controlled experiments show that the improvement is not simply due to a larger number of parameters or a specific choice of fusion module. The ablation study also shows that removing or replacing individual GCM components reduces the model's performance.

The current study is limited to a single hazard risk prediction task and uses a static set of geographic priors. Future work will focus on evaluating the proposed framework across multiple hazard types and larger geographic regions to better assess its generalisation capability. Incorporating temporal satellite imagery along with updated GIS data could also enable the model to reflect changes in geographic environments over time. In addition, dynamic factors such as traffic patterns, weather conditions, and population movement may be integrated to further improve hazard risk estimation.

Future research will also investigate how the generated risk maps can be incorporated into disaster response and operational safety applications, to help assess the practical use of the predicted risk information in real-world systems.

\bibliographystyle{ACM-Reference-Format}
\bibliography{references}

\end{document}
